% c00 —— 题名页（由 main.tex 生成）+ 摘要 + §1 镜像对称（p001 / 原书 p.127）
%
% [原书 p.127] 题名页：ミラー対称性と K3 曲面 / Mirror symmetry and K3 surfaces
%   大阪大学理学研究科 植田一石（Kazushi Ueda, Graduate School of Science, Osaka University）。
%   译本标题页由 main.tex 的 \maketitle 生成，此处不重复。
%   页底刊头「数理解析研究所講究録 第 1918 巻 2014 年 127-151」与页码 127 不排。
%
% 原书 p.127 末句「これは Gromov-Witten 不変量を用いて…」在页边界处截断；
% 整句归 c01（从「这是把 Gromov–Witten 不变量用于……」起），本块止于上一句句号。

\vspace{0.5em}

镜像对称是起源于超弦理论的数学现象，指的是一个空间的复几何与另一个空间的辛几何之间
存在着种种奇妙的关系。对 K3 曲面而言，由于 Hodge 数没有自由度、且亏格 $0$ 的
Gromov–Witten 不变量为零，乍看之下拓扑镜像对称与古典镜像对称似乎平凡；
但实际上它有着丰富的内容。本文希望沿着 [HNU] 来介绍这一点：该文从镜像对称的观点，
详细研究了与 $\mathbb{Q}(\sqrt{5})$ 的对称 Hilbert 模曲面相伴的 K3 曲面族的周期。

\section{镜像对称}

关于镜像对称的猜想有多种，其中代表性的有
\begin{itemize}[leftmargin=2.2em]
  \item 拓扑镜像对称
  \item 古典镜像对称
  \item 同调镜像对称
  \item Strominger–Yau–Zaslow 猜想
\end{itemize}
等。拓扑镜像对称断言：对任意 $n$ 维 Calabi–Yau 流形 $Y$，都存在另一个 $n$ 维
Calabi–Yau 流形 $\check{Y}$，使得 Hodge 数之间成立
\begin{equation}
  h^{i,j}(Y)=h^{i,n-j}(\check{Y}) \tag{1.1}
\end{equation}
这样的关系。古典镜像对称则预言 $Y$ 的 Gromov–Witten 不变量的母函数与 $\check{Y}$ 的
周期之间的奇妙关系。
